% !TeX root = ../main.tex
\chapter{Geometry from the ground up: building a seam}\label{app:D}\markboth{Appendix D}{}
Two strips of paper come with the instruction ``identify these edges.'' What does this instruction give us? At first, only a rule specifying which points are to count as one. It has not told us how to differentiate across the join, how to compare two vectors attached to the joined point, or how to transport the value of a field around the completed strip. We will construct each of these operations. At the end two flat strips will form a spatial circle. Different, locally indistinguishable ways of identifying their field values will give different allowed wave frequencies.

This appendix can be read before Chapter~\ref{ch:01}. It assumes sets and functions, elementary calculus of several real variables, and linear algebra, but not a course in topology or bundles. A useful question to keep asking is: \emph{what is the domain and what is the codomain of the map I am using?} A map of point sets, its derivative, a map of fibers over it, and the pullback of a differential form go between different spaces.

\section{Give the set of points a topology}\label{D-topology}

A \emph{topology} on a set $X$ is a collection of subsets called \emph{open}: it contains $\varnothing$ and $X$, is closed under arbitrary unions, and is closed under finite intersections. A set is \emph{closed} if its complement is open. A neighborhood of $p$ contains an open set containing $p$. In $\mathbb R^d$, $U$ is open when each $p\in U$ is the center of some ball contained in $U$. What we retain in passing to a topological space is the possibility of testing a statement near each point, not the notion of distance that supplied these particular balls.

A \emph{basis} is a collection of open sets such that every open set is a union of members. Euclidean balls with rational centers and positive rational radii suffice: inside any ball about $p$, one can find a smaller ball of this sort that still contains $p$. A space with a countable basis is called \emph{second-countable}. A different condition, \emph{Hausdorff separation}, says that any two distinct points admit disjoint open neighborhoods. Euclidean space has this property: take balls of radii less than half the distance between the points. These two conditions are part of the definition of a manifold in this book; looking Euclidean near each point does not by itself guarantee them globally.

A map $f:X\to Y$ is \emph{continuous} if the inverse image $f^{-1}(V)=\{x:f(x)\in V\}$ of every open $V\subset Y$ is open in $X$. Here $f^{-1}(V)$ makes sense even if $f$ has no inverse function. It suffices to test a basis of $Y$, since inverse images preserve unions. A \emph{homeomorphism} is a bijection continuous in both directions. A \emph{cover} is a family of sets whose union is the space. A space is \emph{connected} when two nonempty disjoint open sets cannot cover it. In the discrete topology every subset is open. A continuous function from a connected strip to the discrete set $\{+1,-1\}$ is therefore constant: the two inverse images of the singletons would otherwise separate the strip.

We need three elementary ways of making spaces. A subset $A\subset X$ has the \emph{subspace topology}: its open sets are $A\cap U$ with $U$ open in $X$. Thus $[0,a)$, for $0<a<L$, is open \emph{in} $[0,L]$. It is not open in $\mathbb R$. The \emph{product topology} on $X\times Y$ has basis $U\times V$ for $U,V$ open. The \emph{disjoint union} $X=R_-\sqcup R_+$ keeps a tag on each point, even if both pieces use the same numerical coordinates; a set in it is open precisely when its part on each piece is open. We use the product to thicken a face, and the disjoint union to start with two genuinely independent regions.

\subsection{How can we tell what is open after identification?}
An \emph{equivalence relation} is reflexive ($p\sim p$), symmetric ($p\sim r$ implies $r\sim p$), and transitive ($p\sim r$, $r\sim z$ imply $p\sim z$). For a list of face pairings, take the relation they \emph{generate}: a point is equivalent to every point reachable by finitely many pairings and their inverses. Its classes form the set $M=X/{\sim}$. Write $q:X\to M$ for the map $p\mapsto[p]$. To specify the \emph{quotient topology}, declare $U\subset M$ open if and only if $q^{-1}(U)$ is open in $X$. This is a topology: take inverse images in the axioms above. The \emph{saturation} of $A\subset X$ is $\operatorname{Sat}(A)=q^{-1}(q(A))$, the union of all classes it meets. The set $A$ is \emph{saturated} when $\operatorname{Sat}(A)=A$. Saturated open sets upstairs are exactly inverse images of open sets downstairs. The image of an arbitrary open set need not be open.

\begin{proposition}[The test for descending a continuous map]\label{D-quotient-test}
Suppose a map $f:X\to Y$ has the same value at equivalent points. There is a unique map $\bar f:M\to Y$ with $f=\bar f\circ q$. It is continuous exactly when $f$ is continuous.
\end{proposition}
\begin{proof}
Define $\bar f([p])=f(p)$. Equality on a class is exactly what makes the formula single valued, and every class has a representative. For an open set $V\subset Y$,
\[
 q^{-1}\big(\bar f^{-1}(V)\big)=f^{-1}(V).
\]
By the quotient topology the set inside $q^{-1}$ is open exactly when the right side is. Applying this to all $V$ proves the test.
\end{proof}

For example, identify $0$ with $L$ in $[0,L]$. A small open arc around the new point has inverse image $[0,a)\cup(L-a,L]$ for $0<a<L/2$. Keeping only the first interval fails: its saturation must contain $L$ but then contains no relative neighborhood of $L$. Figure~\ref{D-arc-figure} shows the missing half of the neighborhood.

\begin{figure}[htbp]\centering
\begin{tikzpicture}[x=.94cm,y=.94cm,>=Stealth,font=\small]
\draw[quiet] (0,0)--(3.9,0);\draw[line width=2pt,marine] (0,0)--(.80,0) (3.10,0)--(3.90,0);
\fill[marine] (0,0) circle (2pt) (3.90,0) circle (2pt);
\draw[marine,fill=white] (.80,0) circle (2pt) (3.10,0) circle (2pt);
\node[below] at (0,-.07) {$0$};\node[below] at (.80,-.07) {$a$};\node[below] at (3.10,-.07) {$L-a$};\node[below] at (3.90,-.07) {$L$};
\draw[->,ink] (4.5,0)--node[above] {$q$} (5.65,0);
\begin{scope}[shift={(7.1,0)}]
\draw[quiet] (0,0) circle (.83);\draw[marine,line width=2pt] (142:.83) arc (142:218:.83);
\draw[marine,fill=white] (142:.83) circle (2pt) (218:.83) circle (2pt);
\fill[marine] (180:.83) circle (2pt);\node[above left,inner sep=1pt] at (-.91,.60) {$[0]=[L]$};
\draw[thin,quiet] (-1.03,.52)--(-.85,.10);
\end{scope}
\end{tikzpicture}
\caption{One open arc downstairs has two relatively open end intervals upstairs. Filled points are included; the two filled points on the interval represent the single marked point on the circle.}\label{D-arc-figure}
\end{figure}

To identify the quotient with the unit circle, use the map
\[
 c:[0,L]\longrightarrow S^1\subset\mathbb R^2,\qquad
 x\longmapsto\big(\cos(2\pi x/L),\sin(2\pi x/L)\big).
\]
Its values agree at the two ends, and no other distinct pair has the same image. It induces a continuous bijection onto the unit circle in $\mathbb R^2$. Why is its inverse continuous? A space is \emph{compact} if every cover by open sets has a finite subcover. The closed bounded interval is compact (the Heine--Borel theorem); its continuous image is compact since one can pull a cover back. A compact subset of a Hausdorff space is closed: separate an outside point from each point of the subset by disjoint neighborhoods, keep finitely many neighborhoods covering the compact subset, and intersect their partners around the outside point. A closed subset of a compact space is compact (add its open complement to a cover). Our continuous bijection thus sends closed sets to closed sets, which is the continuity test for its inverse. This proof uses the topology on the quotient, not how round a circle can be drawn.

\begin{exercise}[Where separation can fail]
Take two copies of the quadrant $(r,y)\in[0,\infty)^2$. On the faces $r=0$ pair equal $y$ only when $y>0$, leaving the two origins distinct. Can a neighborhood of one origin in the quotient be disjoint from a neighborhood of the other?
\end{exercise}
\noindent\emph{Solution.} Pull both neighborhoods back. Each contains points $(0,y)$ for all sufficiently small positive $y$ on its own copy. Such a pair represents the same point downstairs, so the two neighborhoods must meet. Closing the equivalence relation by also pairing the origins changes the space; it is not a choice of coordinates on the original quotient.

\section{Turn two regional charts into one smooth chart}\label{D-collars}
A \emph{chart} is a homeomorphism from an open neighborhood of a point onto an open subset of a coordinate model. Use $\mathbb R^d$ for an interior and $[0,\infty)^k\times\mathbb R^{d-k}$ at a corner with $k$ meeting walls. Here ``open'' means relatively open in the model. In Euclidean open sets, \emph{smooth} means having continuous partial derivatives of all orders. A function on the model is \emph{smooth up to its walls} if locally it is the restriction of such a function on an open set of $\mathbb R^d$. Smooth maps are defined componentwise this way. A change of charts is the map $\psi\circ\varphi^{-1}$ on the overlapping coordinate images. A collection of charts covering a Hausdorff, second-countable space, whose changes are smooth with smooth inverse, gives a \emph{smooth atlas}; adjoining all compatible charts does not change what is smooth. A bijection of such spaces smooth with smooth inverse is a \emph{diffeomorphism}. A homeomorphism only knows open sets: $x\mapsto x^3$ on the line, for instance, has continuous inverse, but its inverse is not differentiable at zero.

At a point of \emph{depth} $k$ exactly $k$ independent nonnegative wall coordinates vanish. The connected pieces of a fixed depth and incidence type are its \emph{strata}; a face in our regions is a specified embedded closed boundary hypersurface. \emph{Embedded} means it has the subspace topology and locally is a coordinate slice, not a curve self-crossing in the picture. A wall ending on another wall does not thereby create a third hypersurface: in $[0,\infty)^2\times\mathbb R^{d-2}$ the intersection of the two independent walls has dimension $d-2$. The number of independent restrictions is its \emph{codimension}. These labels distinguish artificial walls that disappear from physical walls that remain.

\subsection{Construct the map and its inverse}
Let $F_\pm\subset R_\pm$ be paired faces, with a specified diffeomorphism $f:F_-\to F_+$. For the moment stay away from corners. A \emph{collar} is a diffeomorphism from $[0,\epsilon)\times F$ onto a neighborhood of $F$ taking $(0,p)$ to $p$. Local widths suffice here. The coordinate $r$ measures inward parameter, not proper distance. We supply collars as data; proving a collar-existence theorem for every choice of further structures is not needed. Make the quotient of $R_-\sqcup R_+$ identifying $p$ with $f(p)$ and \emph{no} points at positive collar parameter.

An example with different coordinates on the two walls is a good test of the inverse map. Suppose the face coordinate on the minus side is $y$ and on the plus side is $z$, with $z(f(p))=2y(p)+1$ (dimensionless chart coordinates). Use collar coordinates $(r,y)$ and $(r,z)$, both with $r\geq0$. On the saturated union of corresponding collar boxes define
\begin{equation}
 \Psi([r,y]_-)=(-r,y),\qquad
 \Psi([r,z]_+)=\big(r,(z-1)/2\big).
 \label{D-signed-chart}
\end{equation}
The brackets mean the class of the indicated regional point. Both formulas give $(0,y)$ at the seam. They yield all $(s,y)$ in a small open rectangle about it, exactly once after identification. For the inverse take the class of $(-s,y)_-$ when $s\leq0$, or of $(s,2y+1)_+$ when $s\geq0$. The forward map is continuous by Proposition~\ref{D-quotient-test}. The inverse is continuous by \emph{finite closed pasting}: if continuous maps on finitely many closed subsets covering a space agree on overlaps, the inverse image of any closed set under their union is a finite union of closed sets. Here the two closed subsets are the half-rectangles \emph{relative to the open rectangle}. Thus we have a chart across the seam. Using only the coordinate $r$ on both sides would have lost which half a point came from.

To use many such charts, explicitly compose their maps on overlaps. If, for example, the new tangential coordinate on the minus side is $Y=3y$ and on the plus side is $Y=3(z-1)/2$, then $Y=3y$ on the whole signed rectangle. If a second proposed chart instead uses $S=s$ for $s\leq0$ and $S=2s$ for $s\geq0$, the transition is a homeomorphism but fails the first derivative test. Each collar separately is perfectly smooth; the two proposed signed charts do not belong to a single smooth atlas. A stated map on the face alone would not have detected the failure.

\subsection{Intersections and global checks}
At $k$ transverse artificial walls meeting $l$ retained exterior walls, make one chart on each of the $2^k$ local sectors in the product model. If $r_i$ measures inward distance in coordinates (not metric length), put $s_i=\sigma_i r_i$ with $\sigma_i=\pm1$; keep the same nonnegative exterior coordinates $\rho_a$. The proposed charts must cover \emph{all} representatives of this neighborhood; identifications must join exactly the common subfaces and preserve all the indicated other coordinates there. Their images fill
\[
 (-\epsilon,\epsilon)^k\times[0,\epsilon)^l\times V,
 \qquad V\subset\mathbb R^{d-k-l}\ \text{open}.
\]
The same inverse-by-sectors proof works. In a plane four quadrants can erase both axes; if the horizontal axis is a physical boundary, two quadrants erase the vertical cut and leave a half-plane. Three half-lines glued at their origins instead form a letter Y: deleting the central point of a small connected neighborhood leaves three components, not the one or two allowed in a one-dimensional manifold with boundary. Its product with tangential coordinates is the branched local model one would create by identifying three ordinary half-spaces along one open face. Counting incident sectors, not just assigning signs to boundary integrals, is essential.

There are two more checks in calling the result a manifold. Here is a way of carrying them out for the finite product constructions of Chapter~\ref{ch:01}. Suppose each of finitely many pairings is a homeomorphism between \emph{closed} faces, with inverse included; suppose the sector test bounds the number of representatives of a class by $b$. To find all points equivalent to one point, follow chains of pairings, removing repeated points from each chain. No more than $b-1$ steps are needed. There are only finitely many such words. Each word is a homeomorphism on a closed domain to a closed range: when composing, restrict the first map to the inverse image of the second domain. Hence the saturation of a closed set $C$ upstairs is closed, being the finite union of the images of $C$ intersected with the word domains. Distinct finite classes have disjoint open neighborhoods upstairs by Hausdorff separation (separate each pair of points and take finite intersections). Within any such neighborhood $U$ of a class,
\[
 X\setminus\operatorname{Sat}(X\setminus U)
\]
is an open \emph{saturated} neighborhood of that class. These neighborhoods descend, proving Hausdorff separation downstairs. For countability, pull a cover by our assembled charts back to the disjoint union upstairs, which has a countable basis. For each basis member contained in one of these pullbacks, choose one such chart. This countable selection still covers: at each representative a basis member fits inside a pullback. Each selected chart itself has a countable Euclidean relative basis. Their union gives a basis on the quotient. Neither check assumed that there was already a whole smooth spacetime.

\begin{exercise}[An external corner survives]
Take two copies of $r\geq0$, $\rho\geq0$, with a common tangential coordinate $y\in\mathbb R$. Pair the entire $r=0$ faces with equal $(\rho,y)$; do not pair any open part of $\rho=0$. Write the chart of the class of $(0,0,0)$ and identify its depth before and after sewing.
\end{exercise}
\noindent\emph{Solution.} The signed variable is $s=-r$ or $s=r$ according to the copy. With $(s,\rho,y)$ the local image is open in $\mathbb R\times[0,\infty)\times\mathbb R$; at $(0,0,0)$ only $\rho$ is a boundary coordinate. A depth-two point on each piece has become a depth-one point. Extending both $r$ and $\rho$ to signed variables would introduce points that no piece supplied.

\section{Tangent vectors and pullbacks}\label{D-tangent-maps}

In a signed collar, $+\partial_s$ points out of the piece $s\leq0$. No path starting on its face and proceeding into that piece has this initial velocity. Yet a smooth field on the closed piece has a derivative in the $+\partial_s$ direction. To see why, consider what a tangent vector does to a function.

A \emph{germ} at $p$ is a smooth function defined near $p$, with two functions regarded as the same if they agree on some neighborhood of $p$. A tangent vector $X$ at $p$ is a real linear operation on such germs, with values in $\mathbb R$, satisfying the product rule
\[
 X(hk)=h(p)X(k)+k(p)X(h).
\]
An operation with this rule is called a \emph{derivation at $p$}. It kills constants, since $X(1)=2X(1)$. In a chart $x=(x^1,\ldots,x^d)$, expand a germ as
\[
 h(x)-h(p)=\sum_i (x^i-x^i(p))h_i(x),\qquad h_i(p)=\partial_i h(p).
\]
To obtain $h_i$, integrate $\partial_i h$ along the segment from $x(p)$ to $x$ with parameter running from zero to one, using a smooth extension in a small coordinate box. At a corner the derivatives at $p$ do not depend on the extension: two extensions agree on the nearby interior, and their derivatives agree at $p$ by continuity. Applying the product rule now gives $X(h)=\sum_i X(x^i)\partial_i h(p)$. Conversely, every choice of the $d$ numbers $X(x^i)$ gives a derivation by this formula. The $d$ coordinate derivations $\partial_i|_p$ therefore form a basis of the \emph{tangent space} $T_p R$. In new coordinates $z^a$ the components of this same vector are $X(z^a)=\sum_i(\partial_i z^a)(p)X(x^i)$.

The space $T_p R$ has dimension $d$ even at a corner. By contrast, let $\gamma(\tau)$ be a smooth curve in $R$ for $0\leq\tau<\eta$, with $\gamma(0)=p$. Each active wall coordinate $r_i$ starts at zero and stays nonnegative, so $(r_i\circ\gamma)'(0)\geq0$. These inequalities describe exactly the cone of possible right velocities: for every tuple satisfying them, the straight ray stays in the chart for a short time. A two-sided curve staying in $R$ has derivative zero in each active wall coordinate at $p$, since each has a local minimum. Neither collection of velocities should replace the full tangent space when differentiating across a cut.

For a smooth map $H:R\to M$, the \emph{differential} $\mathrm{d}H_p:T_p R\to T_{H(p)}M$ is defined by $(\mathrm{d}H_p X)(h)=X(h\circ H)$. Applying this to the target coordinate functions gives the matrix of partial derivatives of $H$, its Jacobian. The \emph{cotangent space} $T_p^*R$ consists of linear maps from $T_p R$ to $\mathbb R$. For instance $\mathrm{d}h_p(X)=X(h)$; the coordinate covectors $\mathrm{d}x^i$ are dual to the coordinate vectors because $\mathrm{d}x^i(\partial_j)=\delta^i_j$. A smooth choice of covector at each point is a one-form. To bring a one-form $\beta$ on $M$ back to $R$, set
\[
 {(H^*\beta)}_p(X)=\beta_{H(p)}(\mathrm{d}H_p X).
\]
Thus vectors go forward and covectors pull back. No inverse of $H$ is required; in particular $H$ can be the inclusion of a face.

For example, write a minus collar in a common two-dimensional chart by the map $\Phi(r,y)=(s,z)=(-r,y+\lambda r)$, where $\lambda$ is constant. It gives
\[
 \mathrm{d}\Phi(\partial_r)=-\partial_s+\lambda\partial_z,
 \qquad \Phi^*\mathrm{d}s=-\mathrm{d}r,\qquad
 \Phi^*\mathrm{d}z=\mathrm{d}y+\lambda\mathrm{d}r .
\]
Consequently a one-form $\alpha=A\,\mathrm{d}s+B\,\mathrm{d}z$ has the regional expression
\[
 \Phi^*\alpha=\bigl(-(A\circ\Phi)+\lambda(B\circ\Phi)\bigr)\,\mathrm{d}r
                  +(B\circ\Phi)\,\mathrm{d}y .
\]
On the face all these maps are $y\mapsto z=y$. Pulling back to the face itself gives only $B(0,y)\,\mathrm{d}y$: a vector tangent to $r=0$ pairs to zero with $\mathrm{d}r$. Thus the face map knows how to compare tangential traces, but it cannot tell us the coefficient of $\mathrm{d}r$ in the collar.

\begin{exercise}[Clean is not the same as transverse]
Two embedded submanifolds $A,B$ meet \emph{cleanly} if $A\cap B$ is a submanifold and $T_p(A\cap B)=T_pA\cap T_pB$ at each point of intersection. They meet \emph{transversely} if $T_pA+T_pB=T_pM$ there. Compare two distinct coordinate planes in $\mathbb R^3$, a plane taken twice, and the curves $y=0$, $y=x^2$ in $\mathbb R^2$ at zero.
\end{exercise}
\noindent\emph{Solution.} Two coordinate planes meet in a coordinate line; their tangent spaces intersect in that line and span $\mathbb R^3$. A plane taken twice meets cleanly but the sum is the plane, not $\mathbb R^3$. The two curves meet in a point, whose tangent space is zero, but both tangent lines are the $x$-axis. Their meeting is neither clean nor transverse. Independent defining covectors at a product corner provide more than a drawing of intersecting faces.

\section{Forms and the oriented flux through a face}\label{D-forms-flux}

A differential $k$-form assigns to each point an alternating, multilinear function of $k$ tangent vectors, with coefficients smooth in a chart. A function is a zero-form. The coordinate $k$-forms $\mathrm{d}x^{i_1}\wedge\cdots\wedge\mathrm{d}x^{i_k}$, $i_1<\cdots<i_k$, form a basis at each point: their values on $X_1,\ldots,X_k$ are the determinants of the matrices with entries $\mathrm{d}x^{i_a}(X_b)$. The \emph{wedge product} concatenates lists of covectors and extends by bilinearity; interchanging two covectors changes the sign, so a repeated covector gives zero. For example,
\[
 (\mathrm{d}s\wedge\mathrm{d}z)(V,W)=V^s W^z-V^z W^s,
 \qquad \iota_V(\mathrm{d}s\wedge\mathrm{d}z)=V^s\mathrm{d}z-V^z\mathrm{d}s .
\]
The second operation, \emph{contraction}, puts $V$ in the first slot and leaves the other slots free. Pullback of a form puts $\mathrm{d}H_p X_i$ in \emph{each} slot, so pullback preserves wedge products.

The \emph{exterior derivative} differentiates coefficients and puts the new covector first. Writing $\mathrm{d}x^I$ for an increasing wedge as above, the rule is $\mathrm{d}(a_I\,\mathrm{d}x^I)=\mathrm{d}a_I\wedge\mathrm{d}x^I$, extended by addition. Thus $\mathrm{d}\alpha=(\partial_s B-\partial_z A)\,\mathrm{d}s\wedge\mathrm{d}z$ for the one-form just computed. Commuting second partial derivatives occur with opposite wedge signs, so $\mathrm{d}^2=0$. The ordinary product rule gives $\mathrm{d}(\omega\wedge\eta)=\mathrm{d}\omega\wedge\eta + {(-1)}^k\omega\wedge\mathrm{d}\eta$ for $\omega$ of degree $k$. On functions the chain rule reads $H^*\mathrm{d}h=\mathrm{d}(h\circ H)$; applying the wedge and product rules to the coordinate expansion then gives $H^*\mathrm{d}\omega=\mathrm{d}(H^*\omega)$ in every degree.

An orientation chooses positive ordered bases of $T_p R$ consistently from point to point. Let $\epsilon$ be a volume form, a top-degree form positive on these bases. Suppose $F$ is given locally by $\rho=0$, with $\rho\leq0$ inside and $\mathrm{d}\rho\ne0$ on $F$. The tangent space $T_p F$ is the kernel of $\mathrm{d}\rho_p$. Covectors vanishing on this kernel form the \emph{conormal line}; a \emph{coorientation} selects one of its two rays continuously. Here $\nu_p=\mathrm{d}\rho_p$ selects the outward ray. It is a covector on the \emph{ambient} space at $p\in F$, not the pullback $j^*\mathrm{d}\rho$ by inclusion $j:F\hookrightarrow R$; the latter vanishes because $\rho\circ j=0$.

There is also a \emph{normal quotient} $T_pR/T_pF$. Two ambient vectors give the same element of this one-dimensional space if their difference is tangent. The conormal line is its dual: a covector that vanishes on all tangents assigns one value to the entire equivalence class. No metric has so far selected a preferred vector representing that class.

Choose an outward transverse vector $N$, meaning $\nu(N)>0$. The \emph{outward-first} boundary orientation makes $e_1,\ldots,e_{d-1}$ positive on $F$ if $\epsilon(N,e_1,\ldots,e_{d-1})>0$. Any other such vector has the form $cN+T$, with $c>0$ and $T$ tangent to $F$. Its tangential part contributes zero by alternation, so the sign does not change. This convention fixes Stokes' formula for a smooth compactly supported $(d-1)$-form $\omega$ up to the faces:
\[
   \int_R\mathrm{d}\omega=\sum_{F\text{ a boundary face}}\int_F j_F^*\omega .
\]
To check the sign, orient the strip $a\leq s\leq b$, $z\in\mathbb R$, by $\mathrm{d}s\wedge\mathrm{d}z$ and use $\omega=\alpha$ with compact support in $z$. Integrating $\partial_s B-\partial_z A$ over the strip gives $\int_{\mathbb R}(B(b,z)-B(a,z))\,\mathrm{d}z$ by the fundamental theorem of calculus. At $b$ the positive boundary direction is $+\partial_z$; at $a$ it is $-\partial_z$. Splitting the integral at an intermediate $s$ gives the upper-end sign on the minus piece and the lower-end sign on the plus piece.

For a face pairing $f:F_-\to F_+$, write $j_\pm$ for the regional inclusions. An oriented sewing requires $f$ to reverse the induced boundary orientations: the same signed transverse direction is outward on one piece and inward on the other. Changing variables in the second face integral gives
\[
 \int_{F_-}j_-^*\omega_-+\int_{F_+}j_+^*\omega_+
   =\int_{F_-}\bigl(j_-^*\omega_- - f^*(j_+^*\omega_+)\bigr).
\]
In particular, restrictions of a single form on the sewn collar cancel there. It is the tangential pullbacks that enter this calculation, not all the coefficients of the regional forms.

Normalize $N$ by $\nu(N)=1$ and put $\epsilon_\nu=j^*\iota_N\epsilon$. For any vector field $V$, the difference $V-\nu(V)N$ is tangent on $F$, so $j^*\iota_V\epsilon=\nu(V)\epsilon_\nu$. If instead one integrates against a positive measure, the \emph{outward flux density} is $\nu(V)|\epsilon_\nu|$. The absolute value means that the measure transforms by the absolute Jacobian, without an orientation sign. Rescaling $\nu$ by a positive function rescales the normalized transverse vector by the inverse factor, up to a tangent vector. Hence the density is unchanged. When adding outward densities one must not insert a second boundary sign (Appendix~\ref{app:A}).

\begin{exercise}[Flux through a moving face]
In the $(t,x)$ plane orient by $\mathrm{d}t\wedge\mathrm{d}x$ and keep $x\leq h(t)$. Find the outward flux of $V$ per unit $|\mathrm{d}t|$ on the curve $j(t)=(t,h(t))$.

\emph{Answer.} Here $\nu=\mathrm{d}x-h'(t)\mathrm{d}t$ and $N=\partial_x$ has $\nu(N)=1$, giving $\epsilon_\nu=-\mathrm{d}t$. Direct pullback gives $j^*(V^t\mathrm{d}x-V^x\mathrm{d}t)=(h'V^t-V^x)\mathrm{d}t$, with coefficients evaluated on the curve. Relative to the boundary orientation $-\mathrm{d}t$, the outward coefficient is therefore $V^x-h'V^t$. In particular a vector moving with the face has zero flux even if $V^x\ne0$.
\end{exercise}

\section{Lorentzian cones and Cauchy surfaces}\label{D-causal-geometry}

A Lorentzian metric $g$ assigns to each tangent space a symmetric bilinear form varying smoothly with the point. In a suitable basis its matrix is $\operatorname{diag}(-1,1,\ldots,1)$. It is therefore nondegenerate: the only vector pairing to zero with every vector is zero. A nonzero $V$ is timelike, null, or spacelike according as $g(V,V)$ is negative, zero, or positive. Timelike and null vectors are called causal. The timelike vectors form two cones; a \emph{time orientation} selects one continuously as future. Given a future timelike vector $T$, a causal $V$ belongs to the future cone or its null boundary exactly when $g(T,V)<0$. Choosing future is independent of orienting the spacetime for integration.

To compare two metrics one must give them the same vectors as arguments: ${(H^*g)}_p(X,Y)=g_{H(p)}(\mathrm{d}H_p X,\mathrm{d}H_p Y)$. For the collar map $\Phi$ above, $\mathrm{d}s^2-\mathrm{d}z^2$ pulls back to $\mathrm{d}r^2-{(\mathrm{d}y+\lambda\mathrm{d}r)}^2$; the squares here are tensor squares, not wedges. For every $\lambda$ the induced metric at $r=0$ is $-\mathrm{d}y^2$, but its value on tangential vectors does not determine the mixed coefficient in the collar. For a smooth metric on the chosen sewing, compare all coefficients and their one-sided derivatives of every order in the \emph{same signed chart}. The two choices of future must also agree there.

The metric identifies vectors with covectors by $V^\flat(W)=g(V,W)$. Its inverse sends $\nu$ to the vector $\nu^\sharp$ determined by $g(\nu^\sharp,W)=\nu(W)$; in coordinates ${(\nu^\sharp)}^\mu=\sum_\lambda g^{\mu\lambda}\nu_\lambda$, where $(g^{\mu\lambda})$ is the inverse metric matrix. For a hypersurface the orthogonal complement of $T_p F$ is thus the line spanned by its raised conormal. Write $g^{-1}(\nu,\nu)=\nu(\nu^\sharp)$. If this is negative, an orthonormal basis starting with the normalized timelike normal leaves only positive directions in $T_p F$: the hypersurface is \emph{spacelike}. If it is positive, the normal removes a positive direction, leaving a timelike direction in $T_p F$: the hypersurface is \emph{timelike}. If it is zero, $\nu^\sharp$ belongs to $\ker\nu=T_p F$ and is orthogonal to all of $T_p F$. The induced metric is then degenerate and the hypersurface is \emph{null}. Its metric normal cannot serve as the transverse vector used in Stokes' formula.

For example take the Minkowski plane with $g=-\mathrm{d}t^2+\mathrm{d}x^2$, future $\partial_t$, and orientation $\mathrm{d}t\wedge\mathrm{d}x$. Set $u=(t-x)/\sqrt{2}$ and $v=(t+x)/\sqrt{2}$. Substituting $\mathrm{d}t=(\mathrm{d}u+\mathrm{d}v)/\sqrt{2}$ and $\mathrm{d}x=(\mathrm{d}v-\mathrm{d}u)/\sqrt{2}$ gives
\[
 g=-\mathrm{d}u\otimes\mathrm{d}v-\mathrm{d}v\otimes\mathrm{d}u
     =-2\,\mathrm{d}u\,\mathrm{d}v,\qquad
 \mathrm{d}t\wedge\mathrm{d}x=\mathrm{d}u\wedge\mathrm{d}v .
\]
Keep $u\leq0$. The inverse metric has $g^{uv}=g^{vu}=-1$ and zero diagonal, so the outward conormal $\mathrm{d}u$ raises to $-\partial_v$, tangent to the cut. Instead $N=\partial_u$ gives the positive boundary form $\mathrm{d}v$. A unit time translation $V=\partial_t$ carries flux $\mathrm{d}u(V)=1/\sqrt{2}$ per unit $|\mathrm{d}v|$. Indeed $\iota_{\partial_t}(\mathrm{d}t\wedge\mathrm{d}x)=\mathrm{d}x$, which pulls back to $\mathrm{d}v/\sqrt{2}$ there.

A causal curve has a consistent time direction: orient it toward the future and its tangent must be future causal. Piecewise smooth curves suffice for the following examples; more generally one can take locally Lipschitz curves with future causal tangent almost everywhere. The \emph{causal future} $J^+(p)$ consists of $p$ and the points reachable from it by such curves. In the null chart $g(V,V)=-2V^u V^v\leq0$ and $g(\partial_t,V)=-(V^u+V^v)/\sqrt{2}<0$ together require $V^u,V^v\geq0$, not both zero. Integrating along the curve, both coordinates are nondecreasing. Conversely a straight segment reaches any point $u,v\geq0$ from the origin, so this entire quadrant is $J^+(0)$ in the full plane.

A spacelike hypersurface is \emph{Cauchy} if every inextendible causal curve meets it exactly once. Inextendible means that the curve cannot be continued within the spacetime at either end, even after reparametrization. In the full plane a future causal curve has increasing $t$ and $|\mathrm{d}x/\mathrm{d}t|\leq1$. If its $t$ range had a finite upper endpoint, the bound would give a limit for $x$ there; the curve could continue along a vertical timelike line. The same argument holds at the lower end, so every line $t=\text{constant}$ is Cauchy. The cut $u=0$ misses all the inextendible null lines $u=u_0\ne0$.

\begin{exercise}[A spacelike curve that misses a light ray]
Consider the entire graph $t=\sqrt{1+x^2}$ in the Minkowski plane. Is it spacelike? Is it Cauchy?

\emph{Answer.} Its tangent for increasing $x$ is $Y=(x/\sqrt{1+x^2})\partial_t+\partial_x$, so $g(Y,Y)=1/(1+x^2)>0$ everywhere. Yet the inextendible null line $t=x$ never meets the graph: $\sqrt{1+x^2}>x$ for all real $x$. Spacelikeness was a test at each point of the graph; the Cauchy condition also tests curves elsewhere in spacetime.
\end{exercise}

For a time-oriented spacetime without boundary, existence of a spacelike Cauchy hypersurface characterizes \emph{global hyperbolicity}. It is not a matching condition on a collar. For example, identifying the time caps $(0,x)\sim (L,x)$ of a flat slab with $L>0$ sews both the metric and the future direction smoothly but closes every vertical timelike line. No surface meets such a repeatedly traversed line exactly once. Even on the globally hyperbolic plane, a single null trace need not determine a wave: both $\phi=0$ and $\phi=u$ solve $-2\partial_u\partial_v\phi=0$ and vanish on $u=0$. A well-posed evolution requires existence, uniqueness, and continuous dependence in the chosen data and solution norms. These tests concern a specified equation and admissible data, including conditions at any surviving physical wall; smoothness of the sewn metric alone does not imply them (Appendix~\ref{app:B}).

\section{Building a bundle from local products}\label{sec:D-bundles}

A tangent vector at $x$ comes with a base point before it comes with any components. A charged field value also belongs to a space over $x$; reporting it as a complex number requires a choice of phase convention there. To compare two reports at a cut, we need to know how each convention names the elements of the fiber.

A smooth \emph{fiber bundle} consists of a total space $E$, a smooth surjection $\pi:E\to M$ (the projection), and local product descriptions with a common model fiber $Y$. More precisely, there is an open cover $\{U_i\}$ of $M$ and there are diffeomorphisms
\[
 \tau_i:\pi^{-1}(U_i)\longrightarrow U_i\times Y,
 \qquad \operatorname{pr}_1(\tau_i(v))=\pi(v) .
\]
Such a map is a \emph{local trivialization}: it identifies each actual fiber $E_x=\pi^{-1}(x)$ over $U_i$ with $Y$. It need not extend to a trivialization over all of $M$. A \emph{section} $s:M\to E$ is a smooth choice of one element of each fiber, so $\pi\circ s=\mathrm{id}_M$. In the $i$th trivialization it is described by a function $z_i:U_i\to Y$ through $\tau_i(s(x))=(x,z_i(x))$. The function $z_i$ is not the field $s$: a second report of the very same value generally uses another function.

On $U_i\cap U_j$ write $\tau_i\tau_j^{-1}(x,w)=(x,t_{ij}(x)w)$, where $t_{ij}(x)$ is a diffeomorphism of $Y$ and juxtaposition denotes its action on $w$. The index on the left is the \emph{destination}: $z_i=t_{ij}z_j$. At a point covered by three patches, passing from $k$ to $j$ and then to $i$ is a change of description of just one element. Hence
\begin{equation}
 t_{ii}=1,\qquad t_{ji}=t_{ij}^{-1},\qquad t_{ij}t_{jk}=t_{ik} .
 \label{eq:D-bundle-cocycle}
\end{equation}
The last equality on triple overlaps is the \emph{cocycle condition}. The rightmost map acts first.

Conversely, start with the products $U_i\times Y$ and smooth transition maps satisfying \eqref{eq:D-bundle-cocycle}. Form their disjoint union, identify $(x,w)_j$ with $(x,t_{ij}(x)w)_i$ on each overlap, and use the quotient topology. Project the class $[x,w]_j$ to $x$. Over $U_i$ there is an explicit chart and inverse:
\[
 [x,w]_j\longmapsto (x,t_{ij}(x)w),
 \qquad (x,v)\longmapsto [x,v]_i .
\]
The first map gives the same answer on a representative from $k$ because $t_{ik}t_{kj}w=t_{ij}w$. It is continuous by the quotient topology, and the inverse is continuous by restriction of the quotient map to the $i$th product. Their chart changes are the prescribed smooth transitions. This constructs a local product rather than merely a collection of sets of the same size. To construct a section of it, give smooth functions $z_i$ with $z_i=t_{ij}z_j$ and set $s(x)=[x,z_i(x)]_i$; the transition rule is precisely what removes the choice of $i$.

In a rank-$n$ real \emph{vector bundle}, $Y=\mathbb R^n$ and the transitions are invertible linear maps. One can add vectors in $E_x$ by adding their columns in any trivialization, since a linear transition gives the same sum in the next one. Complex bundles use $\mathbb C^n$ instead. The \emph{tangent bundle} $TM$ is the collection of all $T_xM$, with projection $(x,V)\mapsto x$. Each coordinate basis gives a local product, and its transition matrix is the Jacobian $J=\partial x_i/\partial x_j$: a column of vector components changes by $V_i=J V_j$. The \emph{cotangent bundle} $T^*M$ is built out of the dual spaces. A column of covector components changes by $\eta_i=(J^{-1})^{\mathsf T}\eta_j$, since the pairing $\eta_i^{\mathsf T}V_i$ must give $\eta_j^{\mathsf T}V_j$ for every $V_j$. A \emph{frame} $p:\mathbb R^n\to E_x$ is the linear isomorphism sending the standard basis to a chosen ordered basis of the fiber. Once $p$ is chosen, every other frame at $x$ is $p g=p\circ g$ for a unique $g\in\mathrm{GL}(n,\mathbb R)$; none is preferred.

Collecting all frames gives a \emph{principal bundle} with group $\mathrm{GL}(n,\mathbb R)$. More generally take a Lie group $G$, a manifold whose group multiplication and inversion are smooth. A principal $G$-bundle $P\to M$ carries a smooth \emph{right} action, $(p g)k=p(gk)$, free and transitive on each fiber: for any $p,q$ in that fiber there is exactly one $g$ with $q=p g$. Its local trivializations identify this action with right multiplication on the $G$ factor of $U_i\times G$. The local section $e_i$ given by $e_i(x)=\tau_i^{-1}(x,1)$ lets us write every point over $U_i$ uniquely as $e_i(x)a_i$. If $e_j=e_i g_{ij}$, then $e_i a_i=e_j a_j$ implies $a_i=g_{ij}a_j$, not $a_j g_{ij}$.

Now let $G$ act smoothly on $W$ from the \emph{left}, so $\rho(gk)w=\rho(g)(\rho(k)w)$. Make the \emph{associated bundle} $P\times_G W$ out of pairs $(p,w)$ by declaring
\[
 [p g,w]=[p,\rho(g)w],\qquad [p,w]\longmapsto\pi(p) .
\]
Equivalently, changing $p$ to $p g$ while holding an element fixed changes its component to $\rho(g^{-1})w$. Thus $[e_j,z_j]=[e_i,\rho(g_{ij})z_j]$: the same principal frames give the component changes for every associated field. With the frame bundle and $W=\mathbb R^n$ in the defining representation, the map $[p,w]\mapsto p(w)$ recovers the original vector bundle. Other left actions need not be linear.

\begin{example}[A line with no global frame]\label{ex:D-sign-line}
Regard the unit circle as the complex numbers of modulus one. Cover it by $U=S^1\setminus\{-1\}$ and $V=S^1\setminus\{1\}$. Their intersection has two components. Glue real lines by $t_{VU}=+1$ on the upper semicircle and $t_{VU}=-1$ on the lower semicircle, using the inverse maps on the return. These smooth transitions give a line bundle by the construction above. A nowhere-zero section would have a constant sign in each of the connected patches. Its two signs would have to be equal on the upper overlap and opposite on the lower one. Thus there is no such section and the bundle cannot be a product line, whose constant unit section is nowhere zero. It still has a zero section.

The transitions preserve $|w|$, so it makes sense to keep only the unit vectors. They give a principal $\{+1,-1\}$-bundle: either unit vector in a fiber determines the other by multiplication by $-1$. Lift a circuit starting at $1\in S^1$ with coordinate $+1$ in $U$. In a continuous lift to these two sheets the coordinate cannot change within a patch. Cross to $V$ on the upper overlap, obtaining $+1$, and back to $U$ on the lower one, obtaining $-1$. The returned sign is the holonomy of this discrete bundle. Switching charts and back \emph{at one point} gives $t_{UV}(x)t_{VU}(x)=1$; the circuit switches at two different points. Its minus sign does not violate the cocycle condition.
\end{example}

\section{Lifting a seam map and comparing connections}\label{sec:D-bundle-lift}

Let a diffeomorphism $f:F_-\to F_+$ identify faces of two regions carrying principal bundles $P_\pm\to R_\pm$. The plus-side fiber to be used at $y\in F_-$ lies over $f(y)$. For any smooth $b:N\to M$ the \emph{pullback bundle} implements this change of base:
\[
 b^*E=\{(y,v)\in N\times E:\ b(y)=\pi(v)\},\qquad (y,v)\longmapsto y .
\]
Over $b^{-1}(U_i)$ use $y$ and the second coordinate of $\tau_i(v)$ as a trivialization. Its transitions are $t_{ij}(b(y))$. A section $s$ pulls back to $y\mapsto(y,s(b(y)))$; inclusion of a face gives ordinary restriction. For instance, pull the twisted line of Example~\ref{ex:D-sign-line} back by the double covering $b(e^{i\phi})=e^{2i\phi}$. Start with a unit vector at $\phi=0$ and continue it by following its sheet. When $\phi$ reaches $2\pi$ the image has circled twice: the returned sign is $(-1)^2=1$. The choices close to a global unit section, smooth because its local components are constant. Multiplying this section by an arbitrary real coordinate trivializes the pullback line.

At a cut, pulling back $P_+|_{F_+}$ by $f$ puts two principal bundles over $F_-$, but does not identify them. Even the identity of the circle cannot lift to a linear isomorphism from the product line to the twisted line: the image of its constant unit section would be nowhere zero.

A seam lift is an equivariant bundle isomorphism $D:P_-|_{F_-}\to P_+|_{F_+}$ with $\pi_+D=f\pi_-$ and $D(p a)=D(p)a$. The first equality is read along the two routes in Figure~\ref{D-lift-diagram}.

\begin{figure}[htbp]\centering
\begin{tikzpicture}[>=Stealth,thick,font=\small]
\node (p) at (0,1.4) {$P_-|_{F_-}$};\node (q) at (4.5,1.4) {$P_+|_{F_+}$};
\node (x) at (0,0) {$F_-$};\node (y) at (4.5,0) {$F_+$};
\draw[->,marine] (p)--node[above] {$D$} (q);\draw[->,ink] (x)--node[below] {$f$} (y);
\draw[->,quiet] (p)--node[left] {$\pi_-$} (x);\draw[->,quiet] (q)--node[right] {$\pi_+$} (y);
\end{tikzpicture}
\caption{A lift maps actual fiber elements; the face map below it maps base points. Starting with a frame over $y$ on the minus side gives base point $f(y)$ along either route. Commutativity specifies where the frame must land, not which frame there to choose.}\label{D-lift-diagram}\end{figure}

Equivalently, $p\mapsto(\pi_-(p),D(p))$ is an isomorphism over the identity of $F_-$ onto $f^*(P_+|_{F_+})$. To construct it, choose frames on corresponding patches of the faces and specify smooth $h_i$ on the minus patches by
\begin{equation}
 D(e_i^-(y)a)=e_i^+(f(y))h_i(y)a,\qquad
 h_i g_{ij}^-=(g_{ij}^+\circ f)h_j .
 \label{eq:D-lift-frames}
\end{equation}
Indeed, $D(e_j^-)$ computed as $D(e_i^-g_{ij}^-)$ has coefficient $h_i g_{ij}^-$ in the plus $i$ frame. Starting with $e_j^+$ instead gives the coefficient $(g_{ij}^+\circ f)h_j$. Equality makes the displayed prescription a single map, whose inverse uses $f^{-1}$ and $h_i(y)^{-1}$. Glue $p$ to $D(p)$, not just $\pi_-(p)$ to $f(\pi_-(p))$. Choose product descriptions of the collar bundles as pullbacks of their face bundles, as in Chapter~\ref{ch:02}, and extend the frames and $h_i(y)$ constantly in the collar parameter. A common group coordinate is $h_i(y)a_-$ on the minus side and $a_+$ on the plus side. Equation~\eqref{eq:D-lift-frames} now makes its transition between patches the same on both halves. The two descriptions agree on the identified face; they do not identify base points at positive collar distance. For associated fields the induced map is $[p,w]\mapsto[D(p),w]$. In the common plus frame the two values to compare on the face are therefore $\rho(h_i(y))z_i^-(y)$ and $z_i^+(f(y))$.

Changing local frames to $e_i^\pm k_i^\pm$ describes the \emph{same} $D$ by $(k_i^+\circ f)^{-1}h_i k_i^-$. In contrast, an allowed re-marking acts on the pieces. With the base fixed, let $u_\pm$ be automorphisms of the whole regional bundles respecting the exterior markings, and use hats for their face restrictions. The pair carries $p\sim D(p)$ to $u_-p\sim u_+D(p)$, so the new identification is $D'=\widehat u_+D\widehat u_-^{-1}$. The \emph{stabilizer} consists of the pairs with $D'=D$. This tests the face identification; an isomorphism of smooth sewn configurations must also be smooth in the sewn charts and carry the specified fields to their counterparts.

Cut the circle example at one point of each overlap and allow either sign at each cut. Denote the first crossing sign from $U$ to $V$ by $a$ and the return sign by $b$. A smooth automorphism of the discrete bundle over either complete arc is a single sign, say $u$ or $v$. The crossings become $vau^{-1}$ and $ubv^{-1}$, so their ordered product on return is $u(ba)u^{-1}=ba$. Taking $u=1$, $v=a^{-1}$ reduces the pair to $(1,ba)$. With no fixed exterior frame there are two classes, each with the diagonal stabilizer $u=v$. Independent signs supplied at the two ends of one arc would falsely identify the classes: they need not extend over that arc. The sewings together with their re-markings form a \emph{groupoid}: its objects are the choices $(a,b)$, and its arrows are the pairs $(u,v)$ carrying one choice to another. Arrows compose by composing the regional maps and every arrow has an inverse. The set of two classes remembers possible values of $ba$, but not the nonidentity arrow fixing each object.

To compare values at \emph{different} points of a path in a vector bundle, choose a connection. Its covariant derivative is a linear operator taking sections to bundle-valued one-forms, with the Leibniz rule $\nabla(\lambda z)=\mathrm{d}\lambda\otimes z+\lambda\nabla z$ for a scalar function $\lambda$. In a frame it therefore has the form $\nabla z=\mathrm{d}z+A z$: apply the derivative to the frame vectors to obtain the columns of $A$. On a tangent vector $X$ its column is $X(z)+A(X)z$. For a matrix group in its defining representation the potential $A$ is a matrix of one-forms valued in the Lie algebra (the tangent space to $G$ at the identity). To construct a connection from these local potentials, require that the resulting derivatives represent the same fiber vector in either frame. Write $z_+=h z_-$ on a common base patch. Then
\begin{align*}
 \mathrm{d}(h z_-)+A_+h z_-
 &=h\,\mathrm{d}z_-+(\mathrm{d}h+A_+h)z_-
 =h(\mathrm{d}z_-+A_-z_-)
\end{align*}
for all columns precisely when
\begin{equation}
 A_+=h A_-h^{-1}-(\mathrm{d}h)h^{-1} .\label{eq:D-bundle-connection}
\end{equation}
The inhomogeneous term cancels the derivative of the frame change. The same calculation in fixed frames gives the action on a connection of an active gauge transformation $z\mapsto h z$. On a face the common-base notation on the plus side means $z_+\circ f$ and $f^*A_+$, with the potentials restricted to their faces: $(f^*A_+)_y(X)=A_{+,f(y)}(df_y X)$. This only compares tangential directions. For a comparison in the full signed collar, use the realized frame change and differentiate it in the collar, not just on the face.

To construct parallel transport, take a path $\gamma:[0,1]\to M$ in a frame patch and solve the matrix initial-value problem
\[
 \dot a(t)=-A(\dot\gamma(t))a(t),\qquad a(0)=1 .
\]
The right side is tangent to $G$ at $a(t)$: it is the Lie algebra element $-A(\dot\gamma(t))$ carried from the identity by right multiplication. The curve $e(\gamma(t))a(t)$ is the transported frame, called the \emph{horizontal lift} of $\gamma$ starting at $e(\gamma(0))$. To start at $e(\gamma(0))c$, use $a(t)c$; transport commutes with the principal right action. Thus an associated element $[e(\gamma(0)),w]$ is carried to $[e(\gamma(t)),\rho(a(t))w]$. On a longer path change frames at each crossing and compose the transports, with the later matrix on the left.

In the defining representation each transported column solves $\dot z=-A(\dot\gamma)z$. If the initial-to-final matrix is $U_\gamma$, covariance of this equation gives the new matrix $h(\gamma(1))U_\gamma h(\gamma(0))^{-1}$. The initial frame change is inverted because the new input column must first be put in the old frame. For a loop the two endpoints coincide, giving conjugation, as in Chapter~\ref{ch:09}. The resulting transport around a loop is the \emph{holonomy}. Even the \emph{product} complex line over $S^1$ with one global unit frame can have any $U(1)$ holonomy. Set $A=-i\alpha\,\mathrm{d}\theta$, $\alpha\in\mathbb R$; the angle is local but its differential agrees on overlaps. Parallel transport gives $z(2\pi)=e^{2\pi i\alpha}z(0)$. Changing components by $h=e^{i n\theta}$, $n\in\mathbb Z$, sends $\alpha$ to $\alpha+n$ without changing this phase. Here the nontrivial phase comes from the connection, not from any twist of the bundle.

A small loop also measures \emph{curvature}. In commuting coordinates $s,y$ define $F_{sy}$ by $[\nabla_s,\nabla_y]z=F_{sy}z$. Expanding the two derivatives cancels the terms containing derivatives of $z$ and leaves
\[
 F_{sy}=\partial_s A_y-\partial_y A_s+[A_s,A_y],\qquad [A_s,A_y]=A_sA_y-A_yA_s .
\]
These are the components of the curvature two-form $F=\mathrm{d}A+A\wedge A$, with matrix multiplication in the wedge product. For arbitrary vector fields $X,Y$ subtract $\nabla_{[X,Y]}$ from the commutator to remove the effect of their own failure to commute. The transformation of $\nabla$ implies $F_+=h F_-h^{-1}$ in a common base chart. For a direct check choose $A=-i B s\,\mathrm{d}y$ with constant $B$ in a product complex line. Around the rectangle $0\leq s\leq a$, $0\leq y\leq b$, oriented by $\mathrm{d}s\wedge\mathrm{d}y$, only the side $s=a$ contributes to $\oint A=-iBab$. The transported phase is $e^{-\oint A}=e^{iBab}$. The earlier circle connection had $F=0$ yet could have nontrivial holonomy: a loop encircling this rectangle bounds a surface in its patch, whereas a circuit generating the circle does not.

\begin{exercise}[Pull back before changing frame]
On a seam interval the plus coordinate is $x=2y$ in terms of the minus coordinate. Take
\[
 h(y)=\begin{pmatrix}1&y\\0&1\end{pmatrix},\qquad
 A_-=\begin{pmatrix}0&0\\1&0\end{pmatrix}\mathrm{d}y .
\]
Find $A_+$ in its own coordinate $x$, and check it on the minus-side parallel column $z_-=(1,-y)^{\mathsf T}$.

\emph{Answer.} The right side of \eqref{eq:D-bundle-connection} is $\left(\begin{smallmatrix}y&-1-y^2\\1&-y\end{smallmatrix}\right)\mathrm{d}y$. This is the \emph{pullback} of the desired potential. Since $\mathrm{d}x=2\,\mathrm{d}y$,
\[
 A_+(x)=\tfrac12\begin{pmatrix}x/2&-1-x^2/4\\1&-x/2\end{pmatrix}\mathrm{d}x .
\]
The plus column pulled back is $h(y)z_-(y)=(1-y^2,-y)^{\mathsf T}$. Its $y$ derivative is $(-2y,-1)^{\mathsf T}$, whereas $(f^*A_+)(\partial_y)$ applied to it gives $(2y,1)^{\mathsf T}$. The column is still parallel, although its ordinary derivative has changed.
\end{exercise}

% !TeX root = ../main.tex
\section{A jet is an equivalence class, not a Taylor series}\label{D-jets}
We now have a base manifold, a field bundle $\pi:E\to M$, and a way to compare fibers. How can a first-derivative Lagrangian take a field as input without depending on which coordinates we used to differentiate it? Fix $p\in M$. Consider all local smooth sections defined near $p$. Declare two sections equivalent to order $k$ if, in one common base-and-fiber chart, their values and all component derivatives of total order up to $k$ agree at $p$. The equivalence class of $u$ is its \emph{$k$-jet} $j^k_pu$. Its base point must be the same to compare it with another jet. The family of these classes is the \emph{jet bundle} $J^kE\to M$. For $k=0$ it is just $E$: only the value survives. Forgetting the derivatives of top order maps $J^kE\to J^{k-1}E$ for $k\geq1$. A Lagrangian on $J^1 E$ assigns a density to first jets, not a field history to test against an equation of motion.

Why does this definition not depend on the coordinates? Take a line as base and let two sections have real coefficients $u(x)=x$ and $v(x)=x+x^2$ in a product chart. At the origin their triples (value, first derivative, second derivative) are $(0,1,0)$ and $(0,1,2)$. In particular they have the same first jet but different second jets. Change both coordinates: use $y=2x$ on the base and $Z=\exp(u)$ in the fiber, so that a section with old coefficient $h$ now has coefficient $Z_h(y)=\exp(h(y/2))$. The chain rule gives
\[
 (Z_h)'=\tfrac12 e^h h',\qquad
 (Z_h)''=\tfrac14 e^h\big((h')^2+h''\big),
\]
where the right sides are evaluated at $y/2$. The new triples are $(1,1/2,1/4)$ and $(1,1/2,3/4)$. Neither a derivative's number nor the origin of the fiber coordinates was invariant. \emph{Equality to a specified order} was. In general an order-$j$ derivative of a composite uses only derivatives of the original section through order $j$ and derivatives of the coordinate changes at the corresponding arguments. Equal old jets therefore give equal new jets; applying the inverse changes proves the converse. This supplies smooth transitions on the jet coordinates (with symmetry in repeated differentiations), not a separate tensor transformation for each derivative.

A \emph{germ} remembers more: two sections give the same germ if they agree on some neighborhood of $p$. Even all jets at one point need not determine a smooth germ. The function zero and the function $e^{-1/x^2}$ for $x>0$, extended by zero for $x\leq0$, have the same jets at $0$ but different germs there. On the positive side every derivative is a polynomial in $1/x$ times the exponential and tends to zero. So a gluing proof based on matching derivatives must use the actual given sections, rather than summing a purported convergent Taylor series.

\subsection{Test two sections for smooth descent}
Work in the signed collar of Section~\ref{D-collars}. First rewrite the two local sections in the \emph{same} base coordinates and the \emph{same} fiber chart, using the specified collar realization of the lift. For instance a one-form $a(r,y)\,dr$ on the minus side has coefficient $-a(-s,y)$ in front of $ds$. For a nonconstant fiber change one must differentiate the change as well as the old field. Only after that operation compare ordinary derivatives. If $u_\pm$ denote the resulting scalar components in their inward variables and the common $y$, the test at the seam is
\begin{equation}
 (-1)^m \partial_y^\alpha\partial_{r_-}^m u_-(0,y)
   =\partial_y^\alpha\partial_{r_+}^m u_+(0,y).
 \label{D-field-test}
\end{equation}
Here $\alpha$ is a list of nonnegative integers, one for each tangential coordinate; $\partial_y^\alpha$ means differentiate that many times in each. For a $C^k$ test impose $m+|\alpha|\leq k$ on components already $C^k$ on closed halves; for smoothness impose all orders on smooth halves. The field value on the seam is its order-zero \emph{trace}. Equality of traces as functions of $y$ supplies all purely tangential equalities but says nothing about $\partial_s u$.

Here is a direct sufficiency proof. At an internal seam the two values paste into a continuous function. Paste the two proposed $s$-derivatives into a continuous function $g$ using the first derivative test. On a coordinate interval crossing zero the fundamental theorem of calculus on the two halves gives
\[
 u(h,y)-u(-a,y)=\int_{-a}^{0}g(s,y)\,ds+\int_0^h g(s,y)\,ds
 \qquad(a,h>0).
\]
The same identity holds on subintervals up to zero, so difference quotients give $\partial_s u(0,y)=g(0,y)$. Along tangential lines at the seam differentiate the common trace. All first partials are now continuous, so $u$ is $C^1$. Apply the same argument to the piecewise first partials, then the second partials, up to the required order. At a finite product junction the coordinate interval splits into finitely many sector intervals, and the same integral argument works for \emph{every mixed partial}. At a remaining exterior wall, smoothness in our convention means local extension across that wall: perform the argument on compatible local extensions and restrict back, as in Chapter~\ref{ch:03}. Necessity in each case is simply restriction of an actual smooth section. This is \emph{descent}: regional representatives satisfying the tests are the restrictions of one object, unique since the quotient points are all accounted for. For a gauge arrow apply the test to its group-valued functions in a common frame, not only to the fields on which it acts.

\begin{exercise}[Two different failures at a face and a vertex]
(i) On $s\leq0$ take $u=s$ and on $s\geq0$ take $u=s+bs^3$. Which jets match at $0$? (ii) On four closed quadrants put $u_\sigma=c_\sigma s_1s_2$. All values on both axes are zero. Which derivative at the common vertex detects differing $c_\sigma$?
\end{exercise}
\noindent\emph{Solution.} (i) Value, first, and second derivatives match. The third derivatives are $0$ and $6b$; if $b\ne0$ the field is $C^2$ but not $C^3$. (ii) Each pure derivative at the vertex is zero, but $\partial_{s_1}\partial_{s_2}u_\sigma=c_\sigma$. The full mixed test detects the discrepancy. These pieces need not even pass first-normal-derivative tests away from the vertex along an open axis; comparing only their restrictions to the axes did not test those derivatives.

\begin{exercise}[Can the metrics be glued in the stated charts?]
Use the inward collars $r_\pm\geq0$ and a common time $t$. Give each half the metric $g_\pm=dr_\pm^2-e^{2a_\pm r_\pm}dt^2$, with real constants $a_\pm$. Both induced metrics on the face are $-dt^2$. For which constants is there a smooth metric in the prescribed signed chart?
\end{exercise}
\noindent\emph{Solution.} The signed coefficients of $dt^2$ are $-e^{-2a_-s}$ and $-e^{2a_+s}$. Equality of the first derivatives at zero requires $2a_-=-2a_+$. In this family that condition is sufficient too: the two exponential formulas are then identical, the $ds^2$ coefficients are both one, and both choose $\partial_t$ as future if so marked. A general smooth metric need not be decided by its first jet. Here the special two-parameter family makes the test exact.

\section{Use a conormal to test an equation}\label{D-symbol}
Field equations can reduce the amount of data that must be exchanged. To see what is geometric about this assertion, let a linear second-order operator in coordinates have highest-derivative part $a^{\mu\nu}(x)\partial_\mu\partial_\nu$. Its \emph{principal symbol} at $p$ and covector $\xi$ is $a^{\mu\nu}(p)\xi_\mu\xi_\nu$ (we omit Fourier factors of $i$). For a system this is a linear map from field components to equation components. The second derivative of a change of coordinates only multiplies first derivatives of the unknown, by the chain rule. The highest-order coefficient therefore has exactly the two Jacobian factors canceled by the transformation of the covector slots. For a nonlinear equation one first linearizes at the solution jet under study. A hypersurface is \emph{characteristic} there if its symbol on a nonzero conormal is not invertible, for a square system of this type. Rescaling that conormal by a nonzero number only multiplies its symbol by the square of that number. This definition uses the equation; calling a hypersurface \emph{null} uses the metric.

In units where the metric light speed is one, put a field of speed $c>0$ in a specified inertial frame on the flat $(t,x)$ plane. Its density is $\frac12(u_t^2-c^2u_x^2-m^2u^2)\,|dt\,dx|$. If $c\ne1$ this is intentionally a propagation law distinct from the metric wave equation. Take the moving cut $s=x-vt=0$, $y=t$. The common coordinate field $U(y,s)=u(y,s+vy)$ has $u_t=U_y-vU_s$, $u_x=U_s$, and the absolute coordinate Jacobian is one. Differentiating the transformed density coefficient with respect to $U_s$ gives the cooriented momentum, and changing coordinates in the equation gives
\begin{align}
 p^s&=(v^2-c^2)U_s-vU_y,\nonumber\\
 0&=U_{yy}-2vU_{ys}+(v^2-c^2)U_{ss}+m^2U .\label{D-moving-symbol}
\end{align}
Suppose two actual half-solutions smooth up to the cut have the same trace $q(y)$ and the same $p^s(y)$. For $v^2\ne c^2$, one first obtains $U_s=(p^s+vq')/(v^2-c^2)$ on both sides. In \eqref{D-moving-symbol} the terms $U_{yy}$ and $U_{ys}$ are now obtained by tangential differentiation. Solving for $U_{ss}$ gives a common second normal derivative. Differentiate the equation in $s$ and repeat to reconstruct every higher normal derivative. Section~\ref{D-jets} turns the matched jets of the half-fields into a smooth solution. A noncharacteristic timelike face need not be a surface on which arbitrary imposed data give a well-posed evolution; this argument started with solutions, not arbitrary data.

For the metric $-dt^2+dx^2$, the squared conormal length is $1-v^2$, but the symbol of our operator on $ds$ is $v^2-c^2$. Thus if $c\ne1$, $v=1$ makes the cut metric-null without making it characteristic. At $v=c$ the momentum instead repeats tangential information: $p^s=-c q'$. For $m=0$ the halves $U_-=0$, $U_+=s$ are solutions with the same zero trace and momentum; their first normal derivatives differ. This failure cannot be repaired by a new choice of transverse coordinate. In Chapter~\ref{ch:07} a nonlinear normal momentum replaces this linear formula, and sharing an invertible branch becomes a separate requirement.

The conormal has one further use later in the book. Collect, for every $p\in\Gamma$, the covectors that vanish on $T_p\Gamma$: this gives the conormal bundle $N^*\Gamma\subset T^*M|_\Gamma$. For a single defining function they are just the multiples of $ds|_p$. For a smooth field the pullback to $\Gamma$ is always defined. For a distribution, the question depends on whether its singular directions include a nonzero conormal; Appendix~\ref{app:C} states the precise sufficient pullback test. That is a third use of $ds$, not another definition of characteristicity.

\section{Pull back a two-form before dividing out a direction}\label{D-reduction}
The preceding tangent spaces belonged to spacetime. On a space of fields a tangent vector $Y$ instead differentiates a family $\phi_\epsilon$ of fields at $\epsilon=0$; on a family of solutions it respects the linearized equation and boundary conditions. We use $\delta$ for exterior differentiation in these variables. A \emph{presymplectic form} is a closed two-form $\omega$ ($\delta\omega=0$); at a point its kernel consists of the $Y$ for which $\omega(Y,Z)=0$ for \emph{every} $Z$. A closed nondegenerate two-form is symplectic in finite dimensions. For the covariant potential of a field the space and the variations have to be specified; a symmetry with a boundary charge need not be in the kernel. We can do a whole reduction with ordinary coordinates to see what the assertion about a kernel means.

In the circle Maxwell example of Chapter~\ref{ch:09}, use arc angles $\alpha_1,\alpha_2$ modulo $2\pi$ for the two open transports $e^{i\alpha_i}$, and real electric momenta $e_1,e_2$. Angle differences by $2\pi$ on local charts are constant, so the one-forms $\delta\alpha_i$ agree globally. The unreduced potential is $\theta=e_1\,\delta\alpha_1+e_2\,\delta\alpha_2$. Impose flux matching first, by the inclusion
\[
 j:(\alpha_1,\alpha_2,e)\longmapsto(\alpha_1,e,\alpha_2,e)
\]
with output ordered as $(\alpha_1,e_1,\alpha_2,e_2)$. Restriction of a form is pullback, so
\[
 j^*\omega=j^*\delta\theta=\delta e\wedge(\delta\alpha_1+\delta\alpha_2).
\]
Contract this with $Y=A\partial_{\alpha_1}+B\partial_{\alpha_2}+C\partial_e$. The result is $C(\delta\alpha_1+\delta\alpha_2)-(A+B)\delta e$. It vanishes exactly for $C=0$ and $A+B=0$. The shifts $(\alpha_1,\alpha_2)\mapsto(\alpha_1+\lambda,\alpha_2-\lambda)$ move along precisely these directions. They are the effective relative endpoint re-marking. The map
\[
 (\alpha_1,\alpha_2,e)\longmapsto(\alpha=\alpha_1+\alpha_2\ \mathrm{mod}\ 2\pi,\ e)
\]
has exactly the shift orbits as fibers; every output has a representative $(\alpha,0,e)$. The resulting cylinder carries $\omega_{\rm red}=\delta e\wedge\delta\alpha$, nondegenerate and pulling back to $j^*\omega$. This is \emph{reduction} in this example: restrict, identify the null directions, and exhibit the form on the quotient. With $\delta H(Y)=\omega(Y,X_H)$ one gets $X_e=\partial_\alpha$, hence $\{e^{i\alpha},e\}=i e^{i\alpha}$. The relative angle was redundant; the sum, the phase of the closed loop, is measurable. The calculation would lose it if both open phases had been discarded at the outset.

\section{Build a field on the circle from two strips}\label{D-capstone}
Here is a problem to try without first guessing a field on an uncut circle. Let $R_i=\mathbb R_t\times[0,d_i]_{x_i}$, $d_i>0$, $i=1,2$. Identify $(t,d_1)_1$ with $(t,0)_2$, and $(t,d_2)_2$ with $(t,0)_1$. Each strip has the flat metric $-dt^2+dx_i^2$ and increasing future time. On each take a real line with its fixed fiber norm, so its orthonormal frame group is $O(1)=\{+1,-1\}$. A field is a section of this line bundle. When going from the end of strip 1 to the start of strip 2, multiply field components by $a\in\{+1,-1\}$; at the return seam multiply by $b\in\{+1,-1\}$. Use these constant identifications throughout the respective collars. On each strip take the massive wave density $\tfrac12[(\partial_t u_i)^2-(\partial_{x_i}u_i)^2-m^2u_i^2]$, $m\geq0$. There is no sheet action. From these data construct the space, identify the field bundle, and find the allowed frequencies of real separated modes. Does the answer depend on which strip end carries a minus sign?

\subsection{Solution: build the base, then the field and its derivative}
Put $L=d_1+d_2$. The circular coordinate is $X=x_1$ on strip 1 and $X=d_1+x_2$ on strip 2, read modulo $L$. The induced map is a homeomorphism onto $\mathbb R\times S^1$; one can write the inverse locally even for a time interval of infinite duration. At the first seam take the signed coordinate $s=x_1-d_1$ on strip 1 and $s=x_2$ on strip 2. At the return seam take $s=x_2-d_2$ on strip 2 and $s=x_1$ on strip 1. The inverse uses the old strip for $s\leq0$ and the new one for $s\geq0$. All transition derivatives for the spatial variable are $+1$ in these oriented charts. The local inverses and direct maps are smooth, both metrics become $-dt^2+ds^2$, and $\partial_t$ agrees. The time slices are Cauchy circles, as in the zero-shift case of Chapter~\ref{ch:01}.

At the first seam a common fiber coordinate is the coefficient on strip 2: it equals $a u_1$ on the old side and $u_2$ on the new. On the covering line, changing the strip-2 frame by $a$ makes the first transition $+1$ and the return transition $H=ba$. Thus the pulled-back component obeys $u(t,X+L)=H u(t,X)$. This is a statement about the chosen field bundle, not an extra force at either seam. The product $H$ is unchanged by admissible constant re-markings of either connected strip. A first derivative of the field in the same common frame also carries the sign $a$, not a second, independently chosen transition.

We can derive the derivative matching without first assuming a smooth field. Vary the regional actions under a common value at the first seam, so $\delta u_2(t,0)=a\,\delta u_1(t,d_1)$. Spatial integration by parts leaves
\[
 \int dt\,\big[-\partial_{x_1}u_1(t,d_1)+a\,\partial_{x_2}u_2(t,0)\big]\,\delta u_1(t,d_1).
\]
Time variations may be compactly supported. Since the face variation is free and $a^2=1$, this gives $\partial_{x_2}u_2(t,0)=a\,\partial_{x_1}u_1(t,d_1)$. The same argument at the second seam gives the return sign $b$. It is the stationarity of this specified action that yields the first derivative, while the fiber identification supplied the common field value. The bulk equation $u_{x_i x_i}=u_{tt}+m^2u$ then recursively supplies all higher spatial derivatives, with time derivatives obtained along the seams. The pieces consequently pass the jet test, not merely a point-value test.

\subsection{Solution: carry two numbers around the assembled space}
For a real mode write $u_i(t,x_i)=\cos(\omega t) f_i(x_i)$. Then $-f_i''=k^2 f_i$ with $k^2=\omega^2-m^2$. The symbol $k^2$ is initially just an eigenvalue; why is it nonnegative? Multiply by $f_i$ and integrate by parts on both intervals. The first seam gives $-f_1(d_1)f'_1(d_1)+f_2(0)f'_2(0)=0$, because both factors acquire $a$ and $a^2=1$. The return seam cancels in the same way. Thus
\[
 k^2\sum_i\int_0^{d_i} f_i^2\,dx_i=\sum_i\int_0^{d_i}(f_i')^2\,dx_i .
\]
For a nonzero mode we may take $k\geq0$. For $k>0$, solving the ordinary initial value problem at either strip's start gives
\[
 Y_i(d)=T_k(d)Y_i(0),\qquad Y_i=\binom{f_i}{f_i'},\qquad
 T_k(d)=\begin{pmatrix}\cos(kd)&\sin(kd)/k\\-k\sin(kd)&\cos(kd)\end{pmatrix}.
\]
For example the two columns are the solutions with initial data $(1,0)$ and $(0,1)$, together with their derivatives. The first join gives $Y_2(0)=a T_k(d_1)Y_1(0)$; the second gives $Y_1(0)=b T_k(d_2)Y_2(0)$. Addition of angles gives $T_k(d_2)T_k(d_1)=T_k(L)$. A nonzero initial vector returns consistently if and only if
\[
 0=\det(\mathbf{1}_2-H T_k(L))=1-H\operatorname{tr}T_k(L)+H^2\det T_k(L)
   =2(1-H\cos(kL)).
\]
Thus for $H=+1$ the positive wave numbers are $2\pi n/L$, $n=1,2,\ldots$; for $H=-1$ they are $(2n+1)\pi/L$, $n=0,1,\ldots$. At these roots $H T_k(L)=\mathbf{1}_2$, the identity on value and slope, so any nonzero initial vector gives a mode with $\omega^2=m^2+k^2$. Check the excluded point $k=0$ separately: $T_0(L)=\left(\begin{smallmatrix}1&L\\0&1\end{smallmatrix}\right)$. Only a constant spatial mode survives for $H=+1$, and none survives for $H=-1$. If both $k=m=0$, the time equation for this one spatial mode also allows a linear function of time. No completeness theorem for mode expansions is needed to decide which separated modes exist.

\begin{exercise}[Hand the data to the next strip]
Let $d_1=d_2=L/2$, $a=-1$, $b=+1$. Start with $Y_1(0)=(1,k)^{\mathsf T}$ at $k=\pi/L$. Find $f_1$ and $f_2$ rather than inferring them from a known antiperiodic solution. What happens if the first join transmits the value but you forget to transmit the slope with the same sign?
\end{exercise}
\noindent\emph{Solution.} The first interval gives $f_1=\cos(kx_1)+\sin(kx_1)$ and $Y_1(d_1)=(1,-k)^{\mathsf T}$. Multiplying both entries by $a$ gives $Y_2(0)=(-1,k)^{\mathsf T}$. Hence $f_2=-\cos(kx_2)+\sin(kx_2)$ and $Y_2(d_2)=(1,k)^{\mathsf T}$, the original data since $b=+1$. Transmitting $(-1,-k)^{\mathsf T}$ instead already violates the stationary interface equation and returns $(-1,k)^{\mathsf T}$ at the final seam. The measurable distinction here is between two products $ba$, not between two positions for the same sign. The local charts, the metric, the incidence of the faces and the regional field equations did not change. What changed was the geometric instruction for comparing a field value after a complete circuit.

